\section{Noise models}

In the last section we considered errors $F$ on a diagram $D$ with edges $E$ to be elements of the displacement group $\hat{\DDD}_{|E|}$.
A more in-depth discussion of errors is warranted; in particular, we will use the notion of the \textit{weight} of an error in order to discuss fault-equivalence in the next section.
In the qubit half of this thesis, when discussing noise models, we swapped to a traditional circuit-based view.
In this qumode half, we will stick with a ZX-based approach.

So, given a diagram $D$ with edges $E$, a \textit{noise model} $\FFF$ will again be a set $\{\FFF_1, \ldots, \FFF_m\}$ of probability distributions, where each $\FFF_j = \{(F_{j,k}, p_{j,k})\}_k$ consists of errors $F_{j,k}$ and their probability $p_{j,k}$ of occurring, and errors $F_{j,k}$ are displacements in $\hat{\DDD}_{|E|}$.
We call errors $F_{j,k}$ that actually appear in one of the probability distributions \textit{atomic errors}.
Atomic errors can be composed to form general errors $F \in \hat{\DDD}_{|E|}$.
The probability of a general error $F$ occurring is the sum of the probabilities of all the different ways it can be made from atomic errors when sampling only one atomic error from each distribution.

The most accurate way to then reason about the behaviour of errors is using these probabilities.
But calculations here can get a bit involved.
We can take a dual approach that's a little more heuristic; rather than assigning each atomic error $F_{j,k}$ a probability $p_{j,k}$ of occurring, we can assign it a \textit{weight} $w_{j,k} = \wt(F_{j,k})$.
For this approach to be useful, a higher weight should correspond to a lower probability.
We can then define the weight of a general error $F \in \hat{\DDD}_{|E|}$ as some combination of weights of atomic errors that can be composed to make $F$.
How exactly this is defined isn't fixed -- we can choose.
If we want to distinguish this weighted approach from the one that used probabilities, we call this a \textit{weighted noise model}.

For example, given diagram $D$ with edges $E$, consider the (weighted) noise model in which atomic errors are displacements $D(\binom{v_j}{u_j})$ on each edge $e_j$, and each such atomic error has weight:
\begin{equation}\label{eq:weight_atomic_error_zx_noise}
    \wt(\hat{D}(\tbinom{v_j}{u_j})) = \twonorm{\tbinom{v_j}{u_j}} = \sqrt{v_j^2 + u_j^2}
\end{equation}
Then define the weight of a general error $\hat{D}(\binom{\bsym{v}}{\bsym{u}})$ as:
\begin{equation}\label{eq:weight_atomic_error_zx_noise}
    \begin{aligned}
        \wt(\hat{D}(\tbinom{\bsym{v}}{\bsym{u}}))
        &= \wt(\hat{D}(\tbinom{v_1}{u_1}) \otimes \ldots \otimes \hat{D}(\tbinom{v_{|E|}}{u_{|E|}})) \\
        &\coloneqq \wt(\hat{D}(\tbinom{v_1}{u_1})) + \ldots + \wt(\hat{D}(\tbinom{v_{|E|}}{u_{|E|}})) \\
        &= \onenorm{\left(\twonorm{\tbinom{v_1}{u_1}}, \ldots, \twonorm{\tbinom{v_{|E|}}{u_{|E|}}}\right)^T}
    \end{aligned}
\end{equation}
That is, we use the 2-norm on each edge, then across different edges we use the 1-norm.
We will call this noise model the \textit{(Focked up) ZX noise model}.

A natural question to ask, perhaps, is why not just use the 2-norm for everything?
That is, for general error $\hat{D}(\tbinom{\bsym{v}}{\bsym{u}})$, why not just define $\wt(\hat{D}(\tbinom{\bsym{v}}{\bsym{u}})) = \twonorm{\tbinom{\bsym{v}}{\bsym{u}})}$?
The intuition behind mixing the 1-norm and 2-norm is to represent noise on each edge being independent, much like the phenomenological noise in the circuit-based approach assumes each qumode experiences noise independently.
